\documentclass[12pt,a4paper]{article}

\makeatletter
\def\input@path{{/home/maenwe/Master_CSM/modèle_cours_latex}}
\makeatother

\usepackage{prise_de_note}
\graphicspath{{images}}
\begin{document}
	\pagenumbering{gobble}
	\begin{titlepage}
		\centering
		
		% Logo
		\begin{figure}[h]
			\centering
			% Ajustez la largeur si besoin (ex: 0.4 ou 0.6)
			\includegraphics[width=0.5\textwidth]{/home/maenwe/Master_CSM/modèle_cours_latex/logo.jpg} 
		\end{figure}
		
		\vspace{2cm}
		
		% Filet horizontal du haut
		\rule{\linewidth}{0.5mm} \\[0.4cm]
		
		% Titre
		{\color{bleuTitre} \Huge \bfseries Phénomène de propagation} \\[0.4cm]
		
		% Filet horizontal du bas
		\rule{\linewidth}{0.5mm}
		
		\vspace{1.5cm}
		
		{\Large \textbf{Fiche de TD 3}}
		
		\vfill
		
		% Bloc Informations
		\begin{minipage}{0.45\textwidth}
			\begin{flushleft} \large
				\textbf{Enseignant :} \\
				M. Benjamin \textsc{BOUTIN}
			\end{flushleft}
		\end{minipage}
		\hfill
		\begin{minipage}{0.45\textwidth}
			\begin{flushright} \large
				\textbf{Étudiant :} \\
				Nathan \textsc{HASCOET}
			\end{flushright}
		\end{minipage}
		
		\vspace{2cm}
		
		% Date
		{\large \today}
		
	\end{titlepage}
	
	\newpage
	\pagenumbering{arabic}
	\setcounter{page}{1}
	
	\begin{sol}
		$ $\\
		\begin{enumerate}
			\begin{minipage}{17cm}
				\item $\Delta u = w(\theta)\frac1r\frac{\partial}{\partial r}\left(rv'(r)\right) + v(r)\frac1{r^2}w''(\theta) = w(\theta)(v''(r)+\frac1r v'(r)) + v(r)\frac1{r^2}w''(\theta)$\\
				
				Ainsi $w(\theta)(v''(r)+\frac1r v'(r)) + v(r)\frac1{r^2}w''(\theta) + (\omega^2-\beta^2)u = 0$.\\
				
				Et finalement : $\frac1{v(r)}(v''(r)+\frac1r v'(r)) = (\beta^2-\omega^2) - \frac1{r^2}\frac{w''(\theta)}{w(\theta)}$.
				
				Ainsi il existe $\lambda \in \R$ tel que : $r^2v''(r)+r v'(r) + (r^2(\omega^2-\lambda^2)-\lambda) v(r) = 0$ et $w''(\theta) = \lambda w(\theta)$.
				
				\item $w(0) = w(2\pi)$, par périodicité, et $w'(0) = w'(2\pi)$. Pour la condition de Dirichlet, $u(R,\theta) = 0$ donc $v(R) = 0$.
				
				Et pour Neumann, $\partial_v u(R,\theta) = 0 = v'(R) w(\theta)$ donc $v'(R) = 0$.
			\end{minipage}
		\end{enumerate}
	\end{sol}
	
\end{document}